\documentclass[onecolumn]{sigchi}

% Week-6 report, restructured 2026-10-29: locked baseline + evaluation
% metrics first, demo results judged against pre-registered thresholds.
\def\plaintitle{CS197 Assignment 6: Progress Report: Numerical Encoding Comparison}
\def\plainauthor{Max Lee, Owen Liu, and Rita Xiang}

\begin{document}
\title{\plaintitle}
\numberofauthors{3}
\author{%
  \alignauthor Max Lee\\
    \email{maxedlee@stanford.edu}\\
  \alignauthor Owen Liu\\
    \email{owen2034@stanford.edu}\\
  \alignauthor Rita Xiang\\
    \email{rita1119@stanford.edu}\\
}
\maketitle

% This Week
\section{This week}
\subsection{Plan}
Last week's plan was to move from evaluating a VAE trained on standardized
numerical features to systematically comparing six numerical encoding
strategies: (1) develop 3--5 testable hypotheses about how feature
distributions respond to different encodings; (2) finalize and debug the
six encoding strategies from Week 2; (3) agree on a fixed VAE configuration
and control for parameter-count differences across encodings; (4) finalize
the evaluation metrics against standardized-feature baselines; and (5)
build a reproducible pipeline and experiment table. All five items are
complete. Following the coaching feedback --- hypotheses and success
criteria first, one variable at a time --- this week we locked the
baseline, the evaluation metrics, and the success thresholds \emph{before}
touching the full data, then dry-ran the whole comparison on the demo data
under those rules (24 parameter-fair runs, 6 encodings $\times$ 3 seeds,
plus fixed-capacity contrasts).

\subsection{Vector}
A comparison of encodings only means something under a locked baseline and
a fixed metric set, so this week we fixed both --- and the demo dry-run
already puts the baseline under pressure. We expected encodings to trade
off prediction against fidelity; instead the two dissociate: the encoding
with the most collapsed latent geometry (mode-specific) gives the best
ranking of the rare outcome, while scalar encodings barely differ and the
binning encoding fails catastrophically. Whether this dissociation is a
property of the encodings or an artifact of the tiny demo sample (11 test
deaths) is the risk the full-data run must retire, judged against the
thresholds locked below.

\section{Evaluation design (locked)}
\subsection{Baseline (locked)}
The primary baseline is \textbf{standard scaling} in the frozen VAE:
standard is the reference for every comparison --- tables report absolute
values with standard first, framed by the ceiling and floor rows, and
every verdict is stated as a difference from standard.
Two reference points frame every table: a \textbf{ceiling}, logistic
regression on the original standardized features $X$ (what compression to
the latent space costs), and a \textbf{chance floor}, outcome prevalence
for AUPRC and 0.5 for AUROC. The baseline is itself a hypothesis under test
(H5), not an assumption: it counts as validated only if it passes the
thresholds below.

\subsection{Evaluation metrics (locked)}
\begin{table*}
\centering
\small
\setlength{\tabcolsep}{4pt}
\resizebox{\textwidth}{!}{%
\begin{tabular}{lll}
Category & Primary metric & Secondary \\
\hline
Downstream prediction (probe on $Z$) & AUPRC & AUROC \\
Feature recovery & Mean per-feature $R^2$ & Per-feature table \\
Geometry preservation & Spearman (300-patient subsample) & Pearson, NN overlap \\
\end{tabular}%
}
\caption{Locked evaluation metrics: three categories, one primary each.
AUPRC is primary because the outcome is rare. Per-encoding diagnostics
(final KL, seed SDs, instability log) are reported but never compared
across encodings; a clustering check (k-means ARI) is exploratory only.}
\end{table*}

\subsection{Protocol and success thresholds (locked)}
Protocol: identical Week 1--2 procedure and frozen 20-feature set on the
full data; stratified splits, seed 42; the six encodings confirmatory;
parameter-fair primary, fixed-width secondary; seeds \{42, 1, 2026\};
latent dimension 8 primary, with 4 and 32 as the secondary axis. Success
thresholds, fixed in advance:
(1) \emph{baseline validated} if standard's geometry $\geq$ 0.80 and no
encoding beats its AUPRC by $\geq$ 0.02;
(2) \emph{dissociation confirmed} if mode-specific beats standard's AUPRC
by $\geq$ 0.02 while trailing its recovery by $\geq$ 0.20;
(3) \emph{binning collapse} if the binning encoding shows recovery
$<$ 0.15 and final KL $<$ 0.5.
Anything that misses its threshold is reported as not confirmed.


\section{Results: Coding and Experiments}
\subsection{Setup and hypotheses}
We used the frozen Week-2 setup throughout: 20 features (12 numerical, 5
binary admission-diagnosis indicators, 3 categorical) from the eICU demo
(1,722 stays), stratified train/validation/test splits (1,101/276/345
stays; 37/9/11 deaths; seed 42), and the Week-3 VAE configuration (latent
dimension 8, 100 epochs, Adam lr $10^{-3}$, batch 256, $\beta = 1$ with a
10-epoch KL warm-up; hidden width 64 in the reference model). Encodings were
compared parameter-fairly: hidden width was adjusted per encoding so total
parameters match the standardized baseline ($\approx$14,700). Mortality at
48 hours is the target $Y$ throughout --- predicted from the latent
representation, never an input feature.

Hypotheses: \textbf{H1} --- skewed features (creatinine, glucose, WBC) are
recovered better under rank-based transforms (quantile-normalized, PLE) than standard scaling.
\textbf{H2} --- mode-specific normalization preserves multimodal structure
best. \textbf{H3} --- binning encodings lose information the VAE cannot
recover. \textbf{H4} --- bounded features (e.g., oxygen saturation) suit
specialized encodings better than plain scaling. \textbf{H5} --- no encoding
transforms the problem: a standardized baseline is already a good
representation.

\subsection{Demo dry-run: results}
Test-set results (mean $\pm$ SD over 3 seeds, parameter-fair), with the
locked ceiling and floor printed as reference rows. Paired per-seed
differences against standard (the same three seeds for every encoding)
qualify the headline: mode-specific's AUPRC margin is $+0.065$ on seed
means but seed-inconsistent (paired $\Delta$: $+0.094$, $+0.106$,
$-0.005$; 95\% CI $-0.087$ to $0.216$), while its recovery deficit is
consistent ($\Delta = -0.420$; 95\% CI $-0.543$ to $-0.296$). Periodic's
AUROC margin is the most seed-consistent probe advantage in the set
($+0.044$; positive on all three seeds; 95\% CI $-0.005$ to $0.093$), and
quantile binning's structural deficits are consistent (recovery
$\Delta = -0.596$, geometry $\Delta = -0.834$). The verdicts below follow
the locked mean-based thresholds; the intervals say how far to trust
them.

\begin{table*}
\centering
\small
\setlength{\tabcolsep}{4pt}
\resizebox{\textwidth}{!}{%
\begin{tabular}{lcccc}
Encoding & Probe AUROC & Probe AUPRC & Recovery $R^2$ & Geometry (Spearman) \\
\hline
Standard & 0.661 $\pm$ 0.020 & 0.072 $\pm$ 0.015 & \textbf{0.649 $\pm$ 0.034} & \textbf{0.890 $\pm$ 0.021} \\
Quantile-normalized & 0.686 $\pm$ 0.015 & 0.074 $\pm$ 0.014 & 0.611 $\pm$ 0.014 & 0.641 $\pm$ 0.032 \\
Mode-specific & 0.669 $\pm$ 0.051 & \textbf{0.137 $\pm$ 0.051} & 0.229 $\pm$ 0.015 & 0.176 $\pm$ 0.012 \\
Quantile binning & 0.522 $\pm$ 0.050 & 0.043 $\pm$ 0.008 & 0.053 $\pm$ 0.015 & 0.056 $\pm$ 0.014 \\
PLE & 0.663 $\pm$ 0.013 & 0.056 $\pm$ 0.001 & 0.345 $\pm$ 0.014 & 0.251 $\pm$ 0.025 \\
Periodic & \textbf{0.705 $\pm$ 0.008} & 0.081 $\pm$ 0.002 & 0.352 $\pm$ 0.006 & 0.374 $\pm$ 0.042 \\
\hline
Ceiling: logistic on $X$ & 0.693 & 0.105 & --- & --- \\
Chance floor & 0.500 & 0.032 & 0.000 & 0.000 \\
\end{tabular}%
}
\caption{Encoding comparison on the test split (345 stays, 11 deaths),
eICU demo data. Validation losses are not comparable across encodings,
since each changes the reconstruction target's scale.}
\end{table*}

\begin{figure}
\centering
\includegraphics[width=\columnwidth]{fig_week6_all_encodings.png}
\caption{Per-metric comparison with error bars across the six encodings,
eICU demo data; bars are grey, with the highest bar in each panel shown in
navy. Findings: no encoding wins overall --- the maximum changes
hands by metric. Standard scaling takes both structural metrics (recovery
0.649, geometry 0.890); mode-specific has the highest mean rare-outcome
ranking (AUPRC 0.137, $\approx$1.9$\times$ standard; seed-inconsistent,
see text) despite the weakest structure, so
prediction and fidelity dissociate; periodic takes the best probe
AUROC (0.705); quantile binning is last on both structural metrics. The
six shown are the team's selected methods: standardized scalar,
quantile-normalized scalar, mode-specific normalization (Gaussian
mixtures), quantile binning, PLE, and periodic.}
\end{figure}

\textbf{Which encoding, for which data.} The six encodings differ only
on the 12 continuous features; the categorical variables are one-hot in
every condition (the fixed baseline agreed in coaching), so no difference
between encodings is attributable to categorical handling, and a
categorical-specific comparison (e.g.\ learned embeddings) is out of
scope. On continuous features, the winner depends on the distribution.
For the most extremely skewed feature (WBC, skew 25), rank-based
transforms win: quantile-normalized recovers it best (0.34, vs standard
$-$0.05), with PLE (0.30) and mode-specific (0.19) also clearing the
$+0.10$ rescue threshold; quantile binning does not ($-$0.01). For
moderately skewed labs (creatinine, glucose, BUN), standard scaling
recovers best (0.78--0.80) --- and quantile normalization actively hurts
creatinine (0.48). For the bounded, ceiling-piled feature (oxygen
saturation), standard and quantile-normalized are best (0.50, 0.49),
while the specialized encodings fail (periodic 0.19, mode-specific 0.09,
quantile binning 0.02). Mode-specific is the choice only when the goal is
ranking the rare outcome and latent fidelity does not matter; periodic
gives the best overall probe AUROC but bottom-half structure. Quantile
binning wins in no circumstance we tested.

\textbf{H1 mostly not supported.} Standard scaling already recovers
creatinine ($R^2$ 0.80), glucose (0.80), and BUN (0.78) best. Only the most
extreme feature, WBC, needs a rank-based transform (quantile-normalized
0.34, PLE 0.30, vs standard $-$0.05).
\textbf{H2 supported with a twist.} Mode-specific normalization has the
worst structure (recovery 0.23, geometry 0.18, KL $\approx$ 0.03) yet the
best probe AUPRC (0.137, $\approx$1.9$\times$ standard). Prediction and
fidelity dissociate. Its AUPRC edge, however, is seed-inconsistent (paired
95\% CI crosses zero), while its structural deficit appears on every seed.
\textbf{H3 supported, more strongly than stated.} Quantile binning
--- the binning method among our six --- collapses in the VAE (recovery
0.05, geometry 0.06, final KL 0.013), under both fairness settings and
all seeds --- a property of the binned representation under this VAE,
not of training.
\textbf{H4 not supported.} Oxygen saturation is recovered best by
scalar encodings (standard 0.50, quantile-normalized 0.49) and worst by
the specialized ones (quantile binning 0.02, mode-specific 0.09, periodic
0.19).
\textbf{H5 mostly supported.} Probe AUROC spans 0.66--0.71 across the
non-binning encodings; the one real challenger is periodic (0.705, ahead
of standard on all three seeds, though its paired CI just crosses zero).
Directly answering last report's open question: the standardized baseline
is a good representation, and no encoding beat it overall --- gains are
metric-specific (mode-specific for ranking the rare outcome, periodic for
AUROC, rank transforms for the most-skewed feature).

\textbf{Verdicts against the locked thresholds (demo).}
\emph{Baseline validated?} No --- standard's geometry passes (0.890
$\geq$ 0.80), but mode-specific beats its AUPRC by 0.065 $\geq$ 0.02, so
under the locked rule the baseline is not validated --- it has a
challenger on seed means on the primary metric.
\emph{Dissociation?} The criterion is met on seed means (AUPRC margin
$+$0.065, recovery deficit 0.42), but it is \emph{not} confirmed
statistically: the AUPRC margin's paired 95\% CI crosses zero, so on the
demo the dissociation is a hypothesis for the full data, not a finding.
\emph{Binning collapse?} Confirmed on demo for quantile binning, the
binning method among the six: recovery 0.053 $<$ 0.15 and final KL 0.013
$<$ 0.5, consistent across all seeds.
Diagnostics: no failed runs in the full grid; in the latent-dimension loop
(4/8/32), standard's recovery rises 0.375 $\rightarrow$ 0.649 $\rightarrow$
0.893, mode-specific's AUPRC edge peaks at dimension 8, and quantile
binning never recovers at any dimension (numerically unstable at 32).
Limitations: this dry-run uses the demo cohort only (1,722 stays, 11
test deaths), a single dataset and a single outcome; no mechanism test has
been run yet, and subgroup performance was not assessed. The full-data
run is the deciding test.

\section{Team members' contributions}
\subsection{Individual Contributions}
This week we locked the evaluation design and ran the systematic
comparison of numerical encoding strategies under it. Owen implemented
the experiment pipeline and the parameter-fair training protocol, ran
the 24-run grid across the six encodings, and produced the comparison
tables and figures.
Max evaluated the six encodings on downstream prediction and feature
recovery; Rita evaluated them on downstream prediction and geometry
preservation. Throughout the week we discussed the
hypotheses, reviewed each other's code and results, and planned the
full-data re-run together.

\end{document}
